% ============================================================
%  R. Nielsen, "Religionsphilosophie."
%  Kjøbenhavn: Forlagt af den Gyldendalske Boghandel (F. Hegel);
%  I. Cohens Bogtrykkeri, 1869.
%  TRANSCRIPTION (Danish) — from the Det Kgl. Bibliotek scan.
%
%  Scan: ~/bibliotek/Nielsen, Rasmus/religionsphilosophie.pdf (557 PDF pp., KB).
%        Copied into this folder as scan.pdf (gitignored).
%  TYPE: Antiqua (Latin type) throughout — NOT Fraktur.
%        Letterspacing (Sperrsatz) -> \emph{}; image-verify per page by zoom.
%        The OCR text layer CANNOT see letterspacing and is badly scrambled
%        ("Prinoiper", "Villi es", "Sønn en's Væ s en"): use it only as a crutch.
%
%  OFFSET (VERIFIED): PDF = printed + 13.
%        printed p.1 = PDF 14; p.17 = PDF 30; p.42 = PDF 55; p.56 = PDF 69;
%        p.82 = PDF 95; p.538 = PDF 551. Title leaf = PDF 8; preface theses = PDF 10.
%
%  EXTENT: body printed pp. 1--537 (KB catalogue: "537 s.").
%
%  STRUCTURE OF THE INDLEDNING (pp. 1--81) — reconstructed from the running text,
%        because the FIRST LEAF OF THE INDHOLD IS MISSING from this scan (the
%        surviving Indhold leaves, PDF 11--13, begin mid-way at "§ 9 ... 129"):
%          Religion og Philosophie                          § 1   pp.  1--10
%          Religion og Mythologie                           § 2   pp. 11--22
%          Aabenbaring og Tro                               § 3   pp. 23--41
%          Ordet og Aanden                                  § 4   pp. 42--55
%          Troesprincipet: Religionsphilosophiens Methode   § 5   pp. 56--81
%        Then the tripartite body: Tro paa Faderen (§ 6, p.82), Tro paa Sønnen
%        (§ 18, p.272), Tro paa Aanden (§ 32, p.458 region).
%
%  ERRATA (from "RETTELSER", printed p.538 / PDF 551) — none fall in pp. 1--81:
%        p.150 Anm.        : unserer      -> unseres
%        p.257 ll.12--13 f.o., l.15 f.o., l.19 f.o., l.20 f.o. — see RESUME-NOTES.md
%        p.284 l.6  f.n.   : forklare     -> bortforklare
%        p.400 l.6  f.o.   : Prophetens   -> Propheternes
%        p.508 l.11 f.n.   : til i sin    -> til sin
%        p.592 l.4  f.o.   : Selvishedeus -> Selvvishedens
%
%  STATUS: IN PROGRESS. Done & image-verified: title leaf, preface theses,
%        and printed pp. 1--42 (Indledning: §§ 1--3 complete -- § 3 "Aabenbaring
%        og Tro", pp.23--41, with run-heads a) Aabenbaring og hellig Historie
%        (p.24), b) Aabenbaring og Fornuft (p.28), c) Aabenbaringen Tilegnelse
%        i Troen (p.35); § 4 "Ordet og Aanden" opens p.42).
%        RESUME AT printed p.43 (PDF 56). See RESUME-NOTES.md.
%        NOTE: the older indledning/transcription.tex is SUPERSEDED by this file.
%        It covered only pp. 1--17 (breaking off mid-page), carried no page-break
%        comments and no Sperrsatz markup, and mis-transcribed the p.17 sentence
%        "absolut uforenelige med al Videnskab" as "absolut ueensartede med
%        Videnskab". Its translation.tex is still live and has not been touched.
% ============================================================
\documentclass[12pt]{book}
\usepackage[utf8]{inputenc}
\usepackage[T1]{fontenc}
% Danish hyphenation where available; falls back so the file still compiles in a
% minimal TeX install (e.g. a sandbox without texlive-lang-european).
\IfFileExists{danish.ldf}{\usepackage[danish]{babel}}{\usepackage[english]{babel}}
\usepackage[protrusion=true,expansion=true]{microtype}
\usepackage[margin=1.4in]{geometry}
\usepackage{setspace}
\usepackage{amsmath}
% Libertinus is the house face; fall back to Latin Modern in a minimal install so
% the file always compiles for a quick structural check.
\IfFileExists{libertinus.sty}%
  {\usepackage{libertinus}\usepackage{libertinust1math}}%
  {\usepackage{lmodern}}
\IfFileExists{textalpha.sty}{\usepackage{textalpha}}{} % polytonic Greek, if it occurs later
\usepackage{fancyhdr}
\usepackage[colorlinks=true,linkcolor=black,urlcolor=black,
  pdftitle={Religionsphilosophie},
  pdfauthor={Rasmus Nielsen}]{hyperref}
\setlength{\headheight}{15pt}
\pagestyle{fancy}
\fancyhf{}
\fancyhead[LE,RO]{\thepage}
\fancyhead[RE]{\textit{Religionsphilosophie}}
\fancyhead[LO]{\textit{\leftmark}}
\setlength{\parindent}{1.5em}
\setlength{\parskip}{0pt}
\onehalfspacing

% --- house macros -------------------------------------------------------
% Centred bold sub-head (e.g. "Religion og Philosophie.") + ToC entry.
\newcommand{\subhead}[1]{%
  \begin{center}{\large\bfseries #1}\end{center}%
  \phantomsection\addcontentsline{toc}{subsection}{#1}\markboth{#1}{#1}}
% Centred "§ N." marker (not in the ToC).
\newcommand{\parmark}[1]{\begin{center}{\bfseries \S\,#1.}\end{center}}
% Centred bold run-head (e.g. "a) Religionsphilosophien som speculativ Videnskab.")
\newcommand{\runhead}[1]{%
  \medskip\begin{center}{\bfseries #1}\end{center}\medskip}
% -------------------------------------------------------------------------

\begin{document}

% ============================================================
% TITLE PAGE — PDF 8
% ============================================================
\begin{titlepage}
\centering
\vspace*{3.0cm}
{\Huge\bfseries Religionsphilosophie\par}
\vspace{1.4cm}
{\large af\par}
\vspace{1.0cm}
{\LARGE\bfseries R.\ Nielsen.\par}
\vspace{3.0cm}
\rule{4cm}{0.4pt}\par
\vspace{2.0cm}
{\large\bfseries Kjøbenhavn.\par}
\vspace{0.4cm}
{\normalsize Forlagt af den Gyldendalske Boghandel (F.\ Hegel).\par}
\vspace{0.3cm}
{\small\emph{I.\ Cohens Bogtrykkeri.}\par}
\vspace{0.6cm}
{\large 1869.\par}
\end{titlepage}

\frontmatter

% ============================================================
% PREFACE THESES — PDF 10 (unnumbered leaf)
% ============================================================
\thispagestyle{empty}
\vspace*{0.10\textheight}
\begin{quote}
\itshape\noindent
Tro og Viden ere som Principer absolut ueensartede.

\medskip

Principernes Ueensartethed gjør intet Brud paa de almeengyldige
Tankelove, men viser Umuligheden af en hvilkensomhelst
videnskabelig Overgang enten fra Viden til Tro eller fra Tro til Viden.

\medskip

Den lange Strid mellem Religion og Videnskab har oprindelig sin
Grund i en uklar Blanding af de ueensartede Principer, og maa
altsaa høre op, naar Blandingen ophører.

\medskip

Den Sætning, at Sandhed i Religionen ikke er Sandhed i Videnskaben,
er ikke længer et Paradox, naar Videns absolute Grændse falder
sammen med det Mysterium, hvorfra Troeserkjendelsen gaaer ud, og
hvortil den vender tilbage.

\medskip

Det er disse Hovedsætninger, som Religionsphilosophien i den
foreliggende Form vil udvikle og begrunde. Om Udviklingen og
Begrundelsen har nogen Gyldighed, maa Prøvelsen afgjøre.
\end{quote}

\vspace{1.2cm}
\begin{flushright}
Kjøbenhavn, den 28de December 1868.\par
\vspace{0.4cm}
R.\ Nielsen.
\end{flushright}

\tableofcontents
\mainmatter

% ============================================================
% BODY BEGINS — printed p.1 (PDF 14)
% ============================================================

% ---- printed p.1 (PDF 14) ----
\begin{center}
  {\Large\bfseries INDLEDNING.}\par\vspace{4pt}
  \rule{2.5cm}{0.4pt}
\end{center}
\phantomsection
\addcontentsline{toc}{section}{Indledning}
\markboth{Indledning}{Indledning}

Religionsphilosophiens Forhold til Religionen kan ikke paa
tilfredsstillende Maade angives i almindelige Begrebsbestemmelser;
dets eiendommelige Natur udkræver en eiendommelig Behandling. Hvad
der siges religiøst: Viis mig din Tro af dine Gjerninger, det maa
siges philosophisk: Viis mig dit Princip af dets Gjennemførelse.

\subhead{Religion og Philosophie.}
\parmark{1}

Er Philosophien i det Hele en Videnskab, saa maa
Religionsphilosophien, saasandt den skal fortjene Navn af
Philosophie, naturligviis ogsaa være en Videnskab, men just det
Særegne, at den er en Videnskab om Religionen, giver den et fra de
øvrige Videnskaber afvigende Præg og paatrykker den et eget Mærke.
Hvad dette betyder, vil bedst kunne forstaaes ved et Blik paa de
Hovedformer, Religionsphilosophien i nyere Tid har antaget.

Den videnskabelige Form maa bestemmes ved Principet.
Grundspørgsmaalet er: have Religion og Philosophie det samme
Princip, eller ere Principerne forskjellige? Dersom Religionens
Princip er forskjelligt fra Philosophiens, hvilket Princip skal
Videnskaben da følge, Philosophiens eller Religionens? Først ved en
Besvarelse af
% ---- printed p.2 (PDF 15) ----
disse Spørgsmaal kan Religionsphilosophiens Begreb og Methode
bestemmes.

Gaaer man ud fra den Forudsætning, at Religion og Philosophie have
samme Princip, saa bliver Religionsphilosophien en \emph{speculativ
Videnskab}; antager man, at Principerne vel ere forskjellige, men at
Sandheden dog maa bedømmes efter Vidensprincipet, saa bliver
Religionsphilosophien en \emph{kritisk Videnskab}. Antager man
endelig Principerne for absolut ueensartede og lægger det religiøse
Princip til Grund ved den videnskabelige Belysning, saa bliver
Religionsphilosophien en \emph{omvendt Videnskab}. Den første af
disse Synsmaader repræsenteres af Hegel; den anden fornemmelig af
L.\ Feuerbach; den tredie er den, som i nærværende Fremstilling
kommer til Anvendelse.

\runhead{a) Religionsphilosophien som speculativ Videnskab.}

Religionens saavelsom Philosophiens Gjenstand er, ifølge Hegel, den
evige Sandhed i dens Objectivitet, Gud og Intet uden Gud og
Explicationen af Gud. Philosophien udvikler kun sig selv, idet den
udvikler Religionen, og idet den udvikler Religionen, udvikler den
sig selv. Saaledes falder Religion og Philosophie da sammen til Eet.
Philosophien er i Virkelighed Gudstjeneste, er Religion, thi den er
Forsagelse af subjective Indfald og Meninger i Beskjæftigelsen med
Gud. Ikke desto mindre er Beskjæftigelsen med Gud i Philosophien en
anden, foregaaer paa en anden Art og Maade end i Religionen. I
Religionen er Gud Bevidsthedens umiddelbare Gjenstand; i
Philosophien er Bevidstheden om Gud Guds egen Bevidsthed om sig
selv: Forskjellen ligger i Formen. Religionen som Religion er
folkelig, er for Alle og just derfor henviist til Forestillingens
billedlige Form; men en saadan Form er uadæqvat for Indholdet.
Formen er endelig, Indholdet uendeligt: dette er et Misforhold. Kan
dette Misforhold hæves; kan det samme Indhold, uden at
% ---- printed p.3 (PDF 16) ----
forandre sig, fremstilles i to aldeles forskjellige Former? Her er
Vanskeligheden, og paa denne Vanskelighed strander Hegel. Thi hvor
det gjælder Sandheden som Sandhed, gaaer Ideen saaledes op i Formen,
at Formen umulig kan forandres, uden at Indholdet forandres. Under
sin Tankebevægelse bliver den speculative Religionsphilosophie
umærkelig ført til et andet Resultat end det, hvorpaa den oprindelig
sigtede. Bevægelsen er denne: det begynder med, at Videnskaben skal
anerkjende og begrunde Religionens Sandhed; Religionens Sandhed
bestemmes ved det Sande i Indholdet, d.\ e.\ ved det større eller
mindre Qvantum Livsidealitet, en given Religion i Forhold til det
Trin af Formudvikling, hvorpaa den befinder sig, er istand til at
udtrykke; da nu Religionens Form ikke er videnskabelig, men den
videnskabelige Form den eneste fuldkomne, saa bringes det til
Bevidsthed, at selv den fuldkomne Religion har en ufuldkommen Form,
og kan paa Grund af Formens Ufuldkommenhed kun udtrykke det Sande
paa ufuldkommen Maade. Det, hvorved den fuldkomne Religion adskiller
sig fra Philosophien, er altsaa en Form, der maa ophæves; men idet
Formen ophæves, ophæves det for Religionen Eiendommelige og dermed
Religionen selv. Ved at miste sin Form mister Religionen tillige sit
Indhold, thi Indholdets Sandhed beroer paa Formen. Altsaa ender det
med, at Philosophien, for at kunne begribe det Sande i Religionen,
maa begribe Religionen selv som en Usandhed. Det negative, fra
Videnskabens Side frastødende Forhold mellem Philosophie og Religion
er med kraftig Conseqvens kommen til Udtalelse i Nyhegelianismen.

\runhead{b) Religionsphilosophien som kritisk Videnskab.}

Medens den speculativt positive Retning nærmest holder sig til det
for Religion og Philosophie fælles Indhold og stræber at udvikle
dette Sandhedens sande
% ---- printed p.4 (PDF 17) ----
Indhold til den sande Form, gaaer den negativt kritiske Retning, i
Erkjendelse af Religionens og Philosophiens principielle
Forskjellighed, nærmest ud fra den for Religionen eiendommelige
Form, paaviser dens øiensynlige Strid med Fornuften, og kommer saa
gjennem en Paaviisning af det Usande i Formen til det Usande i
Indholdet. Umiddelbart seer det ud, som havde Religionen med lutter
virkelige, om end overnaturlige, Kjendsgjerninger at gjøre; men den
til Viden udviklede Fornuft begriber, at de overnaturlige
Kjendsgjerninger ikke ere virkelige, og de virkelige ikke
overnaturlige. Det er ikke Aanden som Aand, men det sjælelige Livs
ideale Naturkræfter, der røre sig oprindeligt i den oprindeligt
religiøse Bevidsthed. De religiøse Forestillinger ere vel ikke
Frembringelser af bevidst Vilkaarlighed, tvertimod de skyldes en
ubevidst Digtning, en Sjælenes Trang til Selvobjectivering, en
psychologisk Nødvendighed, hvormed den af Ideen overvældede
Phantasie har stillet det Guddommelige i Mennesket frem for
Mennesket ved at stille det udenfor Mennesket. Religionen er
menneskelig; den er den menneskelige Fornufts første elementære
Gjennembrud; den har sit Udspring, ikke i Guds Kjærlighed, men i
Menneskets Lidenskaber. Hjertets Drømme, Sindets Rørelser, Sjælens
Længsler indeholde de Grundelementer, af hvilke den ideebeaandede
Phantasie har dannet alle Troens Guder, ogsaa den ene Almægtige.
Dette stadfæstes ved Indblik saavel i Mono\-theismens som i
Polytheismens Væsen. Den mægtige, vrede, hævnfnysende Jehova i det
Gamle Testamente er, videnskabelig seet, en idealiseret Speiling af
Jødefolkets abstracte, mægtige Aand, af Israels stærke, lovbundne og
dog selvraadige Jeg. Christus i det Nye Testamente er ikke et
virkeligt Gudmenneske, thi et virkeligt Gudmenneske er en
Modsigelse, men dog væsentlig en Menneskegud, det saarede Hjertes,
den ængstede Sjæls forsonende Gud, en Gud for de Lidende. For
Phantasie og Følelse fremstiller det Ideale sig umiddel\-
% ---- printed p.5 (PDF 18) ----
bart; derpaa beroer ikke blot Ønsket selv, men ogsaa dets
Opfyldelse. Mennesket med sin indsnevrende Naturbegrændsning ønsker
sig Magt; havde det Magt som det har Villie, vilde det føle sig
saligt og frit; men ak desværre det er afmægtigt; saa føler det i
sin Afmagt Trang til guddommelig Bistand. Den guddommelige Bistand
bestemmes umiddelbart saaledes, at en Gud, der har Magt over
Naturen, vilkaarlig griber ind i Tingenes Gang og hjælper ved et
Mirakel. I Miraklet anskues den ubegrændsede, fra Naturtvang
frigjorte Villie, en Villie, der formaaer at tilintetgjøre al
Modstand, tilfredsstille alle Ønsker og høre alle Bønner. Miraklet
er Troens Grund, Haabets Udvei, Phantasiens Trøst, det fromme Gemyts
hellige Eventyr; naar en Gud er bleven saa gammel og svag, at han
ikke længere kan gjøre Mirakler, gjør han bedst i at resignere,
ophøre at være Gud, og ligesaa maa en Religion, der ikke
vedblivende kan opvise Mirakler, ophøre at være Religion. Her er det
Vendepunkt, hvor Striden mellem Religion og Philosophie culminerer.
Religionen holder paa Miraklet, Philosophien paa Fornuften. Miraklet
er i Modsigelse med Fornuften; naar Fornuften vaagner til
Erkjendelse af Virkelighedens Love, forsvinder Miraklet, og med
Miraklet tillige Religionen. Saaledes har da Philosophien, som
negativ Kritik, ifølge egen Forsikkring, fuldbyrdet Religionens
Opløsning.

Ad den negative Vei er der imidlertid naaet et positivt Resultat.
Resultatet er, at Religionen, forskudt og forhaanet af
„Videnskaben“, bliver henviist til sit eget Princip og derved
opfordret til at stole paa egne Kræfter og hævde sin Selvstændighed
overfor Videnskaben. Det afgjørende Brud mellem Religion og
Videnskab kan da paa disse Vilkaar ligesaavel komme Religionen som
Videnskaben tilgode. Religionens absolute Ueensartethed med al
objectiv Viden og Videnskab er med gjennemtræn\-
% ---- printed p.6 (PDF 19) ----
gende Skarphed og seirrig Overlegenhed paaviist af S.\ Kierkegaard.

Grundspørgsmaalet, hvorfra S.\ Kierkegaard gaaer ud, er dette:
\emph{Hvorvidt kan Sandheden læres?} Her møder saa, hvad Sokrates i
Meno kalder en stridslysten Sætning: at et Menneske umuligt kan
søge, hvad han veed, og ligesaa umuligt kan søge, hvad han ikke
veed, thi hvad han veed, kan han ikke søge, da han veed det, og hvad
han ikke veed, kan han ikke søge, thi han veed jo end ikke, hvad det
er, han skal søge. Sokrates hæver Vanskeligheden derved, at al Læren
og Søgen kun er en Erindren, saa at den Uvidende blot behøver at
mindes, for ved sig selv at besinde sig paa, hvad han veed.
Sandheden bringes saaledes ikke ind i ham, men var i ham. Dette er,
hvad vi kunne kalde den humane Forklaring. Fra dette Synspunkt faaer
„Øieblikket“, det Guddommeliges Indtræden i Tiden, Underet, ingen
Betydning. „For den sokratiske Betragtning er ethvert Menneske sig
selv det Centrale, og hele Verden centraliserer kun til ham, fordi
hans Selverkjendelse er en Gudserkjendelse.“ Den, der lærer
Sandheden, lærer den ved sig selv, Læreren er kun Anledningen;
Øieblikket, da Sandheden indlyser, er i Forhold til den værende,
Menneskets Væsen iboende Sandhed et forsvindende og forsaavidt
ligegyldigt Tidspunkt. „\emph{Udgangspunktet er et Intet}; thi i
samme Øieblik som jeg opdager, at jeg fra Evighed af har vidst
Sandheden uden at vide det, i samme Nu er hiint Øieblik skjult i det
Evige, indoptaget deri saaledes, at jeg, saa at sige, end ikke kan
finde det, selv om jeg søgte det, fordi her intet Her og Der er, men
kun et \textit{ubique et nusquam}.“ Anderledes fra det exclusivt
religiøse Standpunkt. Her forudsættes Mennesket at være, ikke i
Sandheden, men i Usandheden. Skal nu den, der ved Slægtens og ved
egen Skyld har tabt Sandheden, komme til Sandheds Erkjendelse,
nemlig den Sandheds Erkjendelse,
% ---- printed p.7 (PDF 20) ----
der fører til Salighed, da bliver det nødvendigt, at Gud ikke alene
oplyser den Formørkede om Usandheden, men tillige ved selv at
indtræde i Tiden og existere som Menneske giver det i Usandheden
værende Menneske Betingelsen for at naae Sandheden. Ved Guds
Selvmeddelelse i Tiden faaer Øieblikket, det bestemte Tidspunkt, og
dermed Underet, absolut afgjørende Betydning. Underet er Forstandens
Paradox. Ved at fastholde Underet og opstille Paradoxet har
S.\ Kierkegaard sikkret Religionen mod Videnskabens Overgreb og
reist et Bolværk mod den negative Kritiks Indbrud; men saa er han
ogsaa med afgjørende Sigte paa sit „Indledningsproblem, ikke til
Christendommen, men til det at blive en Christen“, bleven staaende
ved den Afgjørelse, at Christendommen ikke er nogen Lære, men en
Existens-Meddelelse, og en Existens-Meddelelse, udtrykt i det
Problem: „at Individets evige Salighed afgjøres i Tiden ved
Forholdet til noget Historisk, der ydermere er saaledes historisk,
at i dets Sammensætning Det er medoptaget, hvilket ifølge sit Væsen
ikke kan blive historisk, og altsaa maa blive det i Kraft af det
Absurde“.

Ifølge den negative Kritik er det Videnskaben, der erklærer
Religionen for et Absurdum og derved mener at opløse den; ifølge den
Kierkegaardske Opgjørelse er det Religionen, der, den objective
Viden til Trods, erklærer sig selv for et Absurdum og derved een
Gang for alle udelukker Videnskaben. Heraf sees, at Religion og
Philosophie ere ikke blot ueensartede, men absolut ueensartede. Er
nu med denne Forudsætning en Religionsphilosophie mulig?

\runhead{c) Religionsphilosophien som omvendt Videnskab.}

Det exclusive, ved Underet bestemte Gudsforhold, der er Religionens
egen uafviselige Forudsætning, kan Videnskaben ikke tilegne sig.
Denne Forudsætning er
% ---- printed p.8 (PDF 21) ----
det Ubegribelige, det Uudgrundelige, og kan kun anerkjendes paa een
Betingelse, den, at den menneskelige Viden anerkjender en Skranke,
en absolut Grændse, hvorved den maa standse. Hvad der i andre
Erkjendelsens Egne viser sig modsigende og urimeligt, er altid et
Lavere, som Videnskaben frit kan hæve sig over; den kan drøfte dette
Lavere, opløse det, og ved dets Opløsning bestandig skride videre
frem. Men Religionen, det for Viden Modsigende, er ikke et Lavere,
men et Høiere, thi Religionen er Livets høieste ideale Magt; hvor
denne Magt i Medfør af sit oprindelige Væsen møder med „et
Absurdum“, der er Veien objectivt ufremkommelig, der finder
Fornuften Gjennemgangen spærret, der maa Videnskaben bøie af.

Men det for den videnskabelige Tænkning Absurde, „Paradoxet“, falder
dog indenfor den tænkende Bevidsthed, saasandt det falder indenfor
Religionen; hvoraf igjen følger, at det for Videnskaben Utænkelige
kan være tænkeligt i Religionen. Der er en væsentlig Forskjel mellem
det: at overskride Videns Grændse og at overskride Tankens Grændse,
mellem det: at falde udenfor Viden og at falde udenfor Tanken.
Forskjellen mellem Tænkeligt og Utænkeligt bliver her at bestemme
ved det Princip, hvorfra Tænkningen gaaer ud. At Paradoxet falder
udenfor Viden, er vist; men falder det derfor udenfor Tanken? For at
Tanken kan støde sig paa Paradoxet, maa Tænkningen jo selv bringe
Paradoxet frem; Paradoxet er kun utænkeligt derved, at det bliver
tænkt; dets dialektiske Knude er en Sammenslyngning af absolut
ueensartede Tanketraade. Existentielt er Knuden uløselig; men deraf
følger ingenlunde, at den ogsaa skulde være det intellectuelt. Den
intellectuelle Løsning er betinget af, at de ueensartede Traade
opredes og skilles fra hverandre. Videnskaben vilde ikke være istand
til at opfatte og bestemme sit Vidensprincip, dersom den ikke
% ---- printed p.9 (PDF 22) ----
ved Modsætningen mellem det, som kan vides, og det, som ikke kan
vides, var istand til, om end blot negativt, at charakterisere det
for Viden uigjennemtrængelige Troesprincip. Men er det Videnskaben
muligt at adskille de absolut ueensartede Principer, da er det den
jo ogsaa muligt at adskille de fra disse Principer udgaaende Tanker,
muligt altsaa at gjøre Forskjel mellem Videnstanker og Troestanker,
mellem Vidensbegreber og Troesbegreber. Med den eiendommelige Opgave
at klare Brydningsforholdet mellem de ueensartede Tanker er
Religionsphilosophien nærmest at bestemme som en
\emph{Grændsevidenskab}.

Religionsphilosophien bliver imidlertid ikke staaende ved
almindelige Grændsebestemmelser for Tro og Viden; den kan umulig
lade Religionen udvikle sit Indhold uden at lade et heelt System af
Troesbegreber komme til Udvikling. Men Troens religiøse Indhold er,
ifølge Principet, et overvidenskabeligt Indhold; her indtræder da
den Modsigelse, at et overvidenskabeligt Indhold optages i
Videnskaben. Dersom Videnskabsformen nu virkelig blev forvexlet med
Indholdets sande oprindelige Form og indført og fastholdt istedenfor
denne, vilde Modsigelsen ganske vist være uopløselig. Men dette er
ingenlunde Tilfældet. Den videnskabelige Reflexion har just den
Bøielighed, at den kan bringe Vidensformen til at reflectere en
anden og modsat Form. Thi ligesom den klare Dag kan bringe Øiet til
at glemme Lysskjæret og see paa Gjenstandene, saaledes kan ogsaa
Reflexionen bringe Tanken til i Viden at see paa Troen og glemme
Viden. Religionsphilosophien er saa langt fra at paatvinge det
religiøse Indhold en fremmed Form, at dens Formudvikling tvertimod
bestaaer i at frigjøre det fra de Tilsætninger af umiddelbart
individuelle Forestillinger, Billeder, Følelser, Villiesyttringer
o.\ s.\ v., hvormed det i Troeslivet selv altid er sammenvoxet.
Forsaavidt Reflexionen, saavel hvad Formen som hvad Indholdet
angaaer, ved Viden ophæver Viden,
% ---- printed p.10 (PDF 23) ----
idet den lægger an paa at klare, ikke en videnskabelig, men en
overvidenskabelig Sandhed, er Forholdet mellem Sandheden og
Videnskaben væsentlig vendt om, og Religionsphilosophien bliver med
Hensyn hertil at charakterisere som en \emph{omvendt Videnskab}.

Det for enhver Videnskab Afgjørende er Begrundelsen. Nu kan jo
vistnok en hvilkensomhelst grundig Indholdsudvikling betragtes som
en eiendommelig Sagudvikling, altsaa som en Udvikling, hvorunder
Sagen udfolder sine Momenter og begrunder sig selv. Men i alle
directe Videnskaber er Sagbegrundelsens Princip i Grunden Eet med
det almene Vidensprincip, saa at Sagbegrundelsen Punkt for Punkt
falder sammen med den paagjældende Videnskabs egen Selvbegrundelse.
I den omvendte Videnskab er Begrundelsen derimod omvendt.
Grundprincipet er her saa langt fra at gaae op i Vidensprincipet, at
de to absolut ueensartede Principer, Underets og Fornuftens, overalt
frastøde hinanden og i deres negative Forhold bestemme hinanden
negativt. Om dette negative Forhold med de dertil svarende
Vexelbestemmelser kan der nu vistnok gives en klar Viden; men den
almene, de modsatte Principer oversvævende Viden er paa nærværende
Trin ikke længer istand til at tage Parti for sit eget Princip; den
omvendte Videnskab har just den Opgave, at sætte Troen i Kraft og
bringe Viden selv til at bøie sig for den. Som det gaaer med det
begrundende Princip, maa det gaae med Begrundelsen. Den omvendte
Videnskab skal hverken fortie eller fornegte de almeengyldige
Vidensgrunde, men den skal bringe det til klar Indsigt, hvorledes
Religionens Selvbekræftelse netop foregaaer derved, at
Vidensgrundene conseqvent opløse sig og give Plads for Troens
absolut selvgyldige Grunde. Ved at begrunde Vidensgrundenes
Omsætning og Opløsning i Troesgrunde begrunder Religionsphilosophien
sin egen videnskabelige Opløsning og faaer paa disse Vilkaar
% ---- printed p.11 (PDF 24) ----
Betydning som en paa Religionen sig offrende, \emph{sig selv i
Troeslære ophævende Videnskab}.

\subhead{Religion og Mythologie.}
\parmark{2}

Af det Forhold, hvori Philosophien sætter sig til Religionen,
afhænger den videnskabelige Opfattelse af Modsætningen mellem sand
og falsk Religion, mellem Religion og Mythologie. Den speculative
Fornuft gjør Mythologien til Religion; den kritiske Forstand gjør
Religionen til Mythologie; kun i den omvendte Viden kan den
principielle Adskillelse mellem Religion og Mythologie gjennemføres.
Vi skulle betragte disse Synsmaader nærmere.

\runhead{a) Mythologie er væsentlig Religion: den speculative\\ Opfattelse.}

Med den Forudsætning, at Sandheden i Religionen og Sandheden i
Videnskaben ikke ere to indbyrdes ueensartede Sandheder, men den ene
og samme absolute Sandhed, maa den speculative Religionsphilosophie
conseqvent anerkjende Religionen i enhver Mythologie, og betragte
den mythiske Indklædning som et Hylster, der let kan borttages.
Religionen er Aand, og for Aanden som Aand ere de endelige
Forestillinger ligegyldige. At udgrunde Aandens Væsen er at udgrunde
Religionens Væsen; for at udgrunde Religionens Væsen udvikler
Philosophien Religionens Begreb. Religionens Begreb er i dets
Almindelighed et Slægtsbegreb, men efter dets særlige
Fremtrædelsesmaader et Artsbegreb. Religionens Idee har sondret sit
Indhold og udlagt sine Momenter i en Fleerhed af mere og mindre
eensidige Religioner, for igjennem disse at samle sig ud af
Adpredelsen og komme
% ---- printed p.12 (PDF 25) ----
til Fuldendelse i den ene, sande, fuldkomne Religion. „Det har“,
siger Hegel, „været Aandens Arbeide gjennem Aartusinder, at udføre
Religionens Begreb og gjøre den til Gjenstand for Bevidstheden.“
Opgaven er at vise, hvorledes Menneske-Aanden, begyndende i al
Naturlighed, har Trin for Trin kunnet hæve sig over den naturbundne
Tilstand til stedse større Frihed og omsider naae det Vendepunkt, da
Aanden begreb sig som Aand og udtrykte Begrebet i den forbilledlige
Existens, der under Samfundets Form skal tilegnes og virkeliggjøres
af Alle. Hegels religionsphilosophiske Systematik er i denne
Henseende charakteristisk. Fra \emph{Naturreligionen}, det Trin,
hvor Bevidstheden i gjærende Eenhed af Natur og Aand anskuer det
Guddommelige, gjennem den \emph{aandige Individualitets eller den
frie Subjectivitets} --- \emph{Høihedens, Skjønhedens og
Hensigtsmæssighedens} --- \emph{Religioner} skrider Udviklingen frem
til Opfattelsen og Bestemmelsen af den i sig klare, aabenbare, til
Aabenbaringsbegrebet svarende Religion, \emph{den absolute Religion,
Christendommen}.

Ifølge denne Synsmaade ere alle Religioner beslægtede, alle
udsprungne af een Rod, alle Grene paa een Stamme. Den gamle
Adskillelse imellem Naturreligionerne og den aabenbarede Religion er
hos Hegel overseet, endog i den Grad, at „Høihedens Religion“ i
Jødedommen bliver --- for Begrebets Skyld --- stillet lavere end
„Skjønhedsreligionen“ hos Grækerne og „Hensigtsmæssighedsreligionen“
hos Romerne. Hvad følger nu heraf? At den jødiske Religion er
Mythologie ligesaavel som den græske og romerske! Men naar
Jødedommen er Mythologie, da er Christendommen ogsaa Mythologie: den
absolute Religion er af samme Oprindelse som de relative. Med den
absolute Religion er det Mythiske altsaa ikke forsvundet; det er
tvertimod blevet Aanden gjennemsigtigt som et til Religionens
eiendommelige Form Med\-
% ---- printed p.13 (PDF 26) ----
hørende. Først naar det i Religionen forestillede Begreb begribes i
Viden, d.\ v.\ s.\ naar Religionen forvandles til Philosophie,
forsvinder det Mythiske.

„Den absolute Religion“, siger Hegel, „er den aabenbare. Religionen
er da først det Aabenbare, naar Religionens Begreb selv er \emph{for
sig}, d.\ v.\ s.\ naar dens Begreb er blevet sig selv objectivt,
ikke i indskrænket, endelig Objectivitet, men saaledes at den efter
sit Begreb er sig objectiv. Vi have saaledes Tvende: Bevidstheden og
Objectet; men i den med sig opfyldte Religion, den aabenbare, der
har fattet sig selv, er Indholdet selv Gjenstanden, og denne
Gjenstand, \emph{det sig vidende Væsen, er Aanden}. Her er først
Aanden som saadan Gjenstand, Religionens Indhold, og Aanden er kun
for Aanden. Idet den er Indhold, Gjenstand, giver den sig selv den
subjective Bevidstheds anden Side, og derved et endeligt Udseende.
Det er Religionen, der er opfyldt med sig selv. Det er denne Idees
abstracte Bestemmelse, eller Religionen er i Virkeligheden Idee. Thi
Idee i philosophisk Betydning er Begrebet, der har sig selv til
Gjenstand, d.\ e.\ har Tilværelse, Realitet, Objectivitet, er ikke
mere det Indre eller Subjective, men objectiverer sig, kun at dets
Objectivitet tillige er dets Tilbagevenden i sig selv, eller,
forsaavidt vi kalde Begrebet Øiemed, det opfyldte, udførte Øiemed,
der følgelig er objectivt.“\footnote{Religionsphilosophie II.
Berlin 1840. S.\ 192--93.}

Men naar Religionens Indhold paa denne Maade falder sammen med
Philosophiens Idee, naar den religiøse Erkjendelse gjøres til Eet
med Speculationens absolute Viden, hvo seer da ikke, endog ved
første Øiekast, Conseqvensen? Religionen fordrer Tro paa sine
overnaturlige Kjendsgjerninger, holder fast ved den existentielle
Umiddelbarhed og vil ikke lade sig omdanne til Philosophie;
Philosophien, der har forvandlet de hellige
% ---- printed p.14 (PDF 27) ----
Kjendsgjerninger til mythiske Forestillinger, er sikker paa Viden,
erklærer al positiv Religion for Mythologie og opgiver Troen.

\runhead{b) Religion er Mythologie: den kritiske Opfattelse.}

Striden mellem Religion og Philosophie begyndte hos Grækerne meget
tidlig, og hvad Romerne angaaer, kunde jo to Augurer paa Ciceros Tid
ikke betragte hinanden uden at lee. Vistnok afgav Religionen det
oprindelige Grundlag for Aandslivet hos Grækerne saavelsom hos
Romerne, men hos dem begge viser det sig ogsaa, at de religiøse
Forestillingers Form kun paa Folkelivets første, barnlige Trin
tilfredsstiller Aanden: Udviklingsloven er, at Bevidstheden med den
tiltagende Reflexion og den fremskridende Cultur voxer fra
Religionen. Da Sokrates havde opdaget Subjectivitetens Dialektik,
sat Selverkjendelsen som Philosophiens Maal og forsaavidt indført
„nye Guder“, kunde Folkereligionen ikke længere modstaae
Philosophernes Angreb. Blev end Mythebegrebet selv ikke i Oldtiden
udviklet med den Klarhed og Dybde, som skeet er i Nutiden, saa blev
det dog, hvad Gudernes Liv og Heroernes Bedrifter angaaer, snart
tilstrækkelig indlysende, at „det Guddommelige“, for at bruge et
Udtryk af Strauss, „ikke kunde være skeet saaledes, at det saaledes
Skete ikke kunde være guddommeligt“. Denne Opløsning, der forklarer
Mytherne, men tilintetgjør Guderne, er ikke et aandeligt
Tilbageskridt, men i Tankens, Fornuftens, Selvbevidsthedens Navn et
væsentligt Fremskridt. Det Mythologiske er det ubevidst Illusoriske;
men dette ligefrem Illusoriske kan Religionen ikke afføre sig uden
at blotte sin medfødte Svaghed og give sig selv til Priis. De
overnaturlige Phænomener i det Nye Testamente ere ligesaa
fornuftstridige som Begivenhederne i den græske Gudekrig. Religion
og Mythologie ere eenstydige Begreber, og Forskjellen mellem lavere
og
% ---- printed p.15 (PDF 28) ----
høiere, ufuldkomnere og fuldkomnere Religioner kun en Forskjel
mellem lavere og høiere, ufuldkomnere og fuldkomnere Mythologier.

Jo høiere en Mythologie staaer, jo mere opfyldt den er af Væsen og
Aand, desto længere vil den naturligviis holde ud med
Culturbevidstheden; men just den Omstændighed, at den aandfuldeste
af alle Myther, Mythen om den Opstandne, har kunnet bære Culturen
gjennem atten Aarhundreder og holde ud med Menneske-Aanden, indtil
Bevidstheden omsider fik Bugt med alle uklare Forudsætninger og kom
til sig selv, idet den „som det frie, for sig værende Almene, blev
sig selv gjennemsigtig“, netop dette viser, hvad det egentlig er,
der i den „aabenbarede“ Religion bliver det Aabenbare. Det er ikke
en Grundforskjel mellem Mythe og Aabenbaring, der kommer for Dagen;
nei, det, der nu kommer for Dagen, er noget ganske Andet, det er den
store objective Sandhed, at Menneske-Aanden, ikke det enkelte
Menneskes Aand, men Menneske-Aanden i sit Princip, er sin egen
Guddom, er i den religiøse Selvfordobling baade Subject og Object,
som endelig Aand Subject, som Aand i indre Uendelighed Object. Dette
er ingen blasphemisk Selvforgudelse, det er jo Opdagelsen af det
Absolute, af Ideen som Idee, af det Subjectives og det Objectives
gudmenneskelige Eenhed.

Ved at gjennemføre Vidensprincipet har den philosophiske Kritik
conseqvent opløst Religionen som Religion --- thi at ville bevare
det Religiøse, naar Religionens Princip er opløst, er naturligviis
Nonsens --- den har da med det Samme naaet det Grændseskjel, hvor
Livsmagterne reagere: det Religiøse gjør sin Ret gjældende, Troen
træder op imod Viden, og de ueensartede Principer maale sig med
hinanden.
% ---- printed p.16 (PDF 29) ----
\runhead{c) Religion er absolut ueensartet med Mythologie: den anti-\\
rationelle Opfattelse.}

Der er saavel i den kritiske som i den speculative Philosophies
Stilling til Religionen en Tvetydighed, der fra Først til Sidst maa
blende Betragteren og forvirre hans Opfattelse, denne Tvetydighed
ligger i Brugen af Ordet Sandhed. At den evige, absolute Sandhed kun
er een, og at den ene guddommelige Sandhed ikke kan være bleven
splidagtig med sig selv eller have deelt sig i to med hinanden
indbyrdes stridende Sandheder, er ganske vist en ligesaa utvetydig
som ubestridelig Forudsætning. Men idet man saa uden videre overseer
Forskjellen imellem en guddommelig og en menneskelig Viden af
Sandheden, overseer, at hvad der i Guds Viden er absolut, maa i
Menneskets være relativt, overseer den principielle Forskjel mellem
religiøs og videnskabelig Erkjendelse, fremkommer Tvetydigheden. Det
seer ud, som om Philosophien virkelig gik ind paa Religionens egne
Forudsætninger, lod disse Forudsætninger udvikle deres religiøse
Conseqvenser, og derved kom til stridige, hinanden ophævende
Resultater. Ja, hvis det virkelig gik saaledes til, da kunde det med
Sandhed siges, at den philosophiske Kritik lod Religionen selv
fuldbyrde sin egen Opløsning. Men saaledes gaaer det ikke til.
Istedenfor at begynde med det religiøse Princip, begynder jo
Videnskaben tvertimod med Vidensprincipet. Ved at anvendes paa
Religionen bliver Vidensprincipet conseqvent et religionsfornegtende
Princip. Naar det da faaer Udseende af, at Kritiken beviser
Religionens Selvopløsning, saa er dette kun et Skin. Sandheden er,
at Kritiken faaer ud, hvad den selv har sat ind; den begynder med at
sætte Religionsfornegtelsen ind, og ender med at faae
Religionsfornegtelsen ud: den opløsende Kritik bevæger sig i en
Cirkel. Thi hvad er det vel, som Videnskaben i Principet sætter som
det Mythiske? Det er just det for
% ---- printed p.17 (PDF 30) ----
Religionen eiendommelige Gudsforhold. At klare og bestemme
Gudsforholdet er at klare og bestemme de for Aandslivet afgjørende
Grundbetingelser. Nu bør dog Videnskaben vide, at disse
Grundbetingelser stille sig anderledes paa Troens end paa Videns
Omraade. Er det virkelig Videnskabens Alvor at faae Aandslivet selv
fra Grunden af opklaret og forstaaet, da maa den, idetmindste indtil
videre, kunne vise den Selvfornegtelse, at see bort fra sit eget
Princip, og lade Religionen komme til Orde.

At Religionsprincipet er absolut ueensartet med Videnskabens
Princip, kan sees af Begrundelsen. Alle Vidensgrunde ere i dybeste
Forstand Nødvendighedsgrunde, Religionens Grunde derimod ere
Villiesgrunde. Guds Villie er Religionens Absolute. Skulde det, hvad
Tankeløse paastaae, være en Modsigelse, at Videnskaben bliver
vidende om et med Viden absolut ueensartet Princip, da maatte det
netop være den negative Kritik, der har hildet sig i Modsigelsen.
Hvad vil denne Kritik? Opløse den positive Religion i Kraft af
Viden. Hvorfor vil den opløse den positive Religion? Fordi dens
oprindelige Forudsætninger ere absolut uforenelige med al Videnskab.
Kan Videnskaben nu ikke blive vidende om de med Videnskaben
uforenelige Forudsætninger, saa behandler Kritiken jo Noget, den
ikke veed hvad er, og dens hele Forehavende er meningsløst;
paastaaer den, at den tilfulde veed, hvad det er, den opløser, da
har den med det Samme indrømmet, at Videnskaben kan blive vidende om
det i Viden Ufordragelige. Men at Religionsprincipet er absolut
ueensartet med Vidensprincipet, vil jo ikke sige Andet end, at dette
i sig eiendommelige Princip, naar det skal opfattes og anvendes som
Vidensprincip, maa blive Videnskaben ufordrageligt. Og hvorfor?
Fordi Religionens Princip er Guds ubetingede Villie.

At kjende Guds Villie, give sig hen i Guds Villie
% ---- printed p.18 (PDF 31) ----
og gjøre Guds Villie, er Religionens Eet og Alt. Men en ubetinget
Villie kan ingen Videnskab assimilere. Vistnok er Villien ikke for
sig uden Væsen og Idee; men derved, at Ideens Indhold sættes i
Villiens Form, og bestemmes af Villien, blive Væsenstankerne
Villiestanker, og Sandheden personlig. Idet Gudsforholdet er sat som
Villiesforhold, er der udtrykkelig sat en Grundforskjel og dermed et
Grændseskjel mellem den religiøse og den videnskabelige, den
personlige og den upersonlige Sandhed. Den upersonlige Sandhed
erkjendes af Ideenødvendighed; den er Menneskets Væsen iboende og
kan ikke tabes, uden at Mennesket taber sin Fornuft og ophører at
være Menneske. Den personlige Sandhed, Villiessandheden, er derimod
identisk med Friheden og kan tabes, naar Friheden tabes, d.\ e.\
naar Gudsforholdet forstyrres. Skal Religionsphilosophien virkelig
begynde med Religionens Princip, da maa den begynde, ikke med en
Viden \emph{af}, men med en Viden \emph{om} to absolut ueensartede
Principer.

Oprindeligt --- saa lyder det i Videnskaben ufordragelige Ord ---
har Mennesket været sat i det begyndende Gudsforhold saaledes, at
det deri har havt alle Forudsætninger og Betingelser for Sandhedens
og Frihedens Erkjendelse. Ved Villiens, altsaa ved Frihedens Misbrug
har det tabt Friheden, er sunken ned i Ufrihed og dermed i personlig
Usandhed. I Usandheden har det tabt Erkjendelsen af Guds Villie, af
Gud som Gud, og derved tillige Betingelsen for at erkjende sig selv
efter sin evige Bestemmelse. Som det faldne Menneskes
Gudserkjendelse er bleven illusorisk, er ogsaa dets Selverkjendelse
bleven illusorisk. Evnen til almindelig Væsenserkjendelse er
bevaret, men Villieserkjendelsen er tabt: Fornuften er bevaret, men
Aanden er formørket. Beroede Gudsforholdet nu derpaa, at Mennesket
ved sig selv af egen aandelig Drift kunde sætte sig i Forhold til
Gud, da vilde det ogsaa beroe paa Mennesket selv, om det efter at
% ---- printed p.19 (PDF 32) ----
være traadt ud af Forholdet paany vilde træde ind og gjenoprette,
hvad det havde forbrudt. Men det er umuligt. Det er ikke Mennesket,
der har begyndt med at sætte sig i Forhold til Gud, men Gud, der fra
Begyndelsen af har sat sig i Forhold til Mennesket; skal det brudte
Forhold gjenoprettes, maa det skee derved, at Gud ved en udtrykkelig
Villiesact paany antager sig Mennesket, altsaa derved, at han
\emph{aabenbarer sig for Mennesket}. Dette er --- vi gjentage det
--- Religionens Grundtanke. Hvad den Vidende fra en oversvævende
Betragtnings objective Standpunkt vil dømme om denne Tanke, om han
vil finde den rimelig eller urimelig, antagelig eller uantagelig,
tiltrækkende eller frastødende, er uvæsentligt; det Væsentlige er,
at Videnskaben bliver vidende om den, mærker sig den og seer,
hvorledes den udvikler sig i Religionen.

Ved den her angivne Grundtanke er Forskjellen mellem Aabenbaring og
Mythe, mellem Religion og Mythologie uafviselig fastsat. Denne
Forskjel bestemmes ikke ved udvortes, men ved indvortes Mærker, ikke
ved historiske Beretninger, men ved Aandens Vidnesbyrd. Men
hvorledes bestemmes saa Aandens Vidnesbyrd? Dersom Aandens
Vidnesbyrd ikke var Andet end de religiøse Individers umiddelbare
Udtalelse af deres egne indre, individuelle Erfaringer, vilde Aanden
være principløs, og Religionen ikke til at adskille fra
Mythologien. Thi hvilken Mangfoldighed af Theophanier, af
Villiesforkyndelser, Raad, Veiledninger, Befalinger, Trusler og
Forjættelser fra Guder til Mennesker indeholdes ikke i de gamle
Gudesagn! Denne Blanding af alskens Indbildninger, denne brogede
Vrimmel af mere og mindre ideale Gudsforestillinger kan ikke blende
det Sind, der holder sig overbeviist om, at Religionen er Aandens
Liv, og at Livet har sit Princip. Det Gudsforhold, der i Eet og Alt
fyldestgjør Aandens, det evige Livs Fordringer, fastsætter Maalet,
% ---- printed p.20 (PDF 33) ----
anviser Veien og tilbyder Midlerne, er og maa være det sande,
ethvert andet derfra afvigende, det usande. Her fordres altsaa en
Kritik, der sikkert adskiller det sande Gudsforhold fra det usande.

Som Principet er, maa Kritiken være. Da nu Principet som
Villiesprincip, det absolute Frihedsprincip, ikke tilhører
Videnskaben, men Religionen, saa følger heraf, at den Kritik, der i
Aand og Sandhed skal kunne skjelne Religion fra Mythologie, maa
være, ikke videnskabelig, men religiøs.

Hvad har saa Videnskaben med denne Kritik at gjøre? Den har at
skaffe de formelle Betingelser tilveie; den har at fremhæve alle
væsentlige Bestemmelser, fjerne alle uvedkommende Hensyn, afværge
alle forvirrende Tankeblandinger, sikkre Bevægelsens conseqvente
Fremgang og blive vidende om det religiøse Princips eiendommelige
Selvudvikling.

Gudsforholdet er et Forhold af Aand til Aand. Herom kan Videnskaben
i al Almindelighed være enig med Religionen; men naar Forholdet
nærmere skal bestemmes, indtræder Adskillelsen. Forholdet har nemlig
to Hovedsider, Ideesiden, der grunder i Fornuftnødvendighed, og den
exclusive Villiesbestemmelse, der grunder i Friheden: hiin er den
abstracte, thi Ideesiden kan fremhæves, uden at Villien erkjendes;
denne er derimod den concrete, thi Villien kan ikke, hverken som
menneskelig eller guddommelig, erkjendes i Sandhed, uden at
erkjendes i Ideen. Saalidt som Aanden er en i sig almindelig Idee
uden Villie, ligesaalidt er den en tom, vilkaarlig Villie uden Idee:
Aandsvillien er Ideevillie, Aanden selv den absolute Eenhed af Idee
og Villie.

I det mythiske Gudsforhold er Friheden lammet, Villieselementet
forqvaklet. Forqvaklingen viser sig saavel fra den guddommelige som
fra den menneskelige Side deri, at Egoismen er traadt i
Aandsvilliens Sted;
% ---- printed p.21 (PDF 34) ----
de mythiske Guder ere Egoister ligesaavel som Menneskene.
Forskjellen er imidlertid den, at kun den menneskelige Egoisme er
reel, den guddommelige opdigtet og phantastisk. Saa vaagner
Fornuften og opdager det Phantastiske; Gudernes Villie forsvinder
med Guderne selv: kun Ideenødvendigheden bliver tilbage. Fra den
aabenbarede Religions Standpunkt kunne Naturreligionerne umulig
antages for at høre til Sandhedens Rige; de ere tvertimod Vidnesbyrd
om et i Usandhed nedsunket Gudsforhold, et forudgaaende Affald fra
den sande Gud. De mythiske Guder ere i Religionens Sprog Afguder,
falske Guder; den Eenhed af Aand og Natur, hvormed Mythologien
begynder, er en falsk Eenhed. Fornuften kan belyse de falske Guder,
opdage Virkelighedens Love og forjage Overtroen; men Ideens
Tvetydighed, den falske Eenhed af Natur og Aand, kan den til sig
selv overladte Fornuft umulig opklare, og det af den gode naturlige
Grund, at Fornuftnødvendighed og Ideenødvendighed væsentlig falde
sammen. Fornuften kan erkjende en Verdensorden i Tingenes
Sammenhæng, kan i Stat, i Kunst, i Videnskab udfolde og klare
Ideernes System; men Aanden som Aand, den evig frie, som absolut
Villie virkende Aand kan den ikke forklare. Aanden som Aand har kun
een Forklaring, og det er dens egen, ubetinget frie
Selvaabenbarelse. Fornuften er saa langt fra at kunne ophæve det
Mythiske, at den ved at miskjende Gudsvillien og lade Friheden paa
tvetydig Maade fødes af Nødvendigheden, meget mere selv i Væsen og
Grund er mythisk. Dette er Religionens Grundtanke; om den vil være
istand til at udvikle og forsvare den imod videnskabelige Angreb,
ville vi faae at see; men Saameget er vist: det er Religionens
Grundtanke, og Videnskaben maa blive vidende derom.

Idet Villiesprincipet er sat som Aandens eget guddommelige Princip,
maa Opfattelsen af det for Religionen
% ---- printed p.22 (PDF 35) ----
Eiendommelige blive, ikke rationel, men antirationel. Det
Antirationelle er ingenlunde det i sig selv Fornuftstridige, men
det, der med en saa eiendommelig Bestemthed staaer Fornuften imod,
at den Fornuftige ved nærmere Eftertanke maa indsee det Ufornuftige
i at ville begribe det. Modsætningen af Religion og Fornuft
bestemmes ingenlunde derved, at Religionen skulde enten gjøre Brud
paa de almindelige Tankelove eller negte Fornuften Ret til at prøve
de objective i Ideenødvendighed begrundede Sandheder; nei,
Modsætningen bestemmes ved det for Religionen eiendommelige Princip,
Villiesprincipet. En videnskabelig Begrundelse af Villiesprincipet
er umulig; men af samme Grund er en videnskabelig Gjendrivelse af
dette Princip da ogsaa umulig. Hvor antirationelt Principet er,
bliver især indlysende, naar vi betænke, at dets Væsen er Underet
selv, at det absolute Villiesprincip netop er Underets Princip.

Naar Talen er om det i Religionen Overnaturlige og „Unaturlige“, som
forarger Fornuften, saa pleier Opmærksomheden gjerne at fæste sig
paa Miraklerne i det Ydre, de udvortes Tegn og Gjerninger, som om
det var dem alene, der vakte Anstød. Herved overseer man det
indvortes overnaturlige Princip. Og naar saa Talen er om dette
Princip, overseer man igjen kun altfor let, at det overnaturlige
Princip er Villiesprincipet selv. Den saakaldte sunde Fornuft finder
det anstødeligt at negte Gud en Villie, men den finder det ikke
anstødeligt at negte Religionen sit Mirakel. Har --- saaledes
raisonnerer man --- Mennesket, som kun er et endeligt Fornuftvæsen,
en, om end ufuldkommen, fri Villie, saa maa Guddommen, det uendelige
Fornuftvæsen, nødvendigviis have en fuldkommen fri Villie. Det
Forargelige sættes altsaa ikke deri, at den Almægtige har en
fornuftig Villie, men deri, at en uendelig fornuftig Villie skulde
kunne beslutte sig til noget saa Ufornuftigt, som at gjøre et
Mirakel. Men det
% ---- printed p.23 (PDF 36) ----
kategoriske Udtryk for den almægtige Villies ubetingede Frihed er
Miraklet. At Gud som Gud kan ville, kan beslutte, kan i en enkelt
bestemt Gjerning udtale en enkelt bestemt Villiestanke, er jo netop
Miraklernes Mirakel, Grundmiraklet: Religionen, Gudsforholdet selv,
er Grundmiraklet.

Forskjellen mellem Religion og Mythologie er da i Korthed denne:
Religionen udtrykker det sande, almægtige, Mythologien det falske,
afmægtige Villiesprincip; Religionen kjendes paa det sande,
Mythologien paa det falske, blot forestillede Mirakel.

\subhead{Aabenbaring og Tro.}
\parmark{3}

Dersom Aandens d.\ v.\ s.\ Villiens Gud ikke var sig selv aabenbar,
vilde Aandsvillien ganske vist heller ikke kunne blive Mennesket
aabenbar; men det for Aabenbaringsbegrebet Afgjørende er dog ikke,
hvad Hegel antager, at Aanden som Aand \emph{bliver} aabenbar, men,
at Aanden selvbestemmende \emph{gjør} sin Villie aabenbar:
Aabenbaringen er som Frihedsact udtrykkelig en Villiesact. Herpaa
beroer fra Guds Side Aabenbaringens uopløselige Eenhed med Underet,
fra Menneskets Side dens uopløselige Eenhed med Troen. Ligesom Troen
bliver guddommelig bestemt ved Aabenbaringen, saaledes bliver
Aabenbaringen igjen menneskelig bestemt ved Troen: Aabenbaring er et
Troesbegreb, Tro et Aabenbaringsbegreb. Den samme Aand, der virker
Aabenbaringen, virker ogsaa Troen; og ligesom Aabenbaringen bliver
aandløs, naar dens objective Elementer adskilles fra Troen, saaledes
bliver ogsaa Troen aandløs, naar den troende Subjectivitet løsriver
sig fra Aabenbaringen. Da begges Eenhed imidlertid ikke er en blot
formel, men en reel Eenhed, altsaa en Modsætningseenhed, vil det til
Sagens
% ---- printed p.24 (PDF 37) ----
Opklaring være nødvendigt at belyse Forholdet fra dets forskjellige
Hovedsider; fra det historiske, ved Modsætningen mellem Aabenbaring
og hellig Historie; fra den metaphysiske, ved Modsætningen mellem
Aabenbaring og Fornuft; fra den psychologiske, ved Aabenbaringens
Tilegnelse i Troen.

\runhead{a) Aabenbaring og hellig Historie.}

Som Villiesact af Gud selv er Aabenbaringen ikke „Manifestation“ af
en speculativ Idee, men guddommelig Handling, en Aandens Gjerning,
en hellig Kjendsgjerning. Aabenbaringen er, ifølge sit Princip,
historisk; kun gjælder det om, at det Historiske ikke kommer ind
under en falsk Belysning. Ved at anlægge samme Maalestok paa den
hellige som paa den profane Historie falder det Kritiken let, ikke
blot at rokke de bibelske Fortællingers Troværdighed, men tillige at
sætte en Skilsmisse mellem Aabenbaring og Tro, der maa tilintetgjøre
begges Betydning. En Betragtning, som hertil svarer, er følgende:
Der viser sig historisk en Forskjel mellem umiddelbar og middelbar
Meddelelse af det, der aabenbares. Den umiddelbare Meddelelse skeer
ikke til Alle, men kun til Enkelte. Disse Enkelte, de Udvalgte, der
saa at sige faae Aabenbaringen paa første Haand, opleve selv det
Vidunderlige. Men hvad et Menneske saaledes oplever, derom har han
jo en ligefrem umiddelbar Vished, men til ligefrem, umiddelbar
Vished svarer Viden, ikke Tro. Jeg troer ikke, at jeg er til, thi
jeg veed det; jeg troer ikke, at jeg virkelig seer og hører, og at
der er synlige og hørlige Ting til, thi jeg veed det. Saaledes med
de Udvalgte; de \emph{troe} ikke, at en Aabenbaring er skeet, thi de
\emph{vide} det. Mængden derimod, de Ikke-Udvalgte, d.\ v.\ s.\ de,
som ikke selv opleve det Vidunderlige, men ved at stole paa
Beretninger om, hvad Andre have oplevet, modtage Aabenbaringen paa
anden Haand,
% ---- printed p.25 (PDF 38) ----
altsaa ret historisk, blive de egentlig Troende. Thi at troe paa en
Aabenbaring, man udtrykkelig har faaet, er jo ingen Sag, men at troe
paa en Aabenbaring, man ikke har faaet, er en anstrengende Opgave.
Paa denne Maade faaer man Aabenbaring og Tro til at falde udenfor
hinanden: \emph{salige ere de, som ikke have seet, men dog troe}.
Ved at falde udenfor Troen bliver Aabenbaringen udvortes objectiv,
altsaa aandløs: ved at falde udenfor Aabenbaringen bliver Troen, med
et tomt Indvortes, hvis Indhold maa søges i et Udvortes, formelt
subjectiv, altsaa aandløs. Denne dobbeltsidige Aandløshed vedkommer
imidlertid ikke Religionen, men alene Kritiken.

For at fyldestgjøre Aandens Fordring, optager og forener den hellige
Historie absolut ueensartede Elementer. Gud forkynder Mennesket sin
Villie; Forkyndelsen skeer umiddelbart, paa endelig bestemt Maade.
Paa den umiddelbare Bestemthed vil det Overnaturlige kjendes; men
den bestemte Umiddelbarhed hefter i Yttringsmaaden tillige ved det
Naturlige, det i Phænomenet Sandselige; Indhold og Form, det Indre
og det Ydre staae i afgjørende Modsætning: Indholdet er
overnaturligt, Formen naturlig, det Indre er oversandseligt, det
Ydre sandseligt. Dette er en Modsigelse. Dersom nu Indholdet var
commensurabelt med Formen, og Forholdet mellem det Indre og det Ydre
et immanent Nødvendighedsforhold, maatte Modsigelsen naturligviis
løses i Viden. Men Indholdet er netop incommensurabelt med Formen;
Forholdet mellem det Indre og det Ydre er ingenlunde begrundet i
Nødvendighed, men tvertimod sat med guddommelig Frihed. Modsigelsen
er, at det Oversandselige skal sandses, det i sig Usynlige sees, det
Uhørlige høres, en saadan Modsigelse maa løses i Troen. Kun Troen i
religiøs Forstand besidder den lykkelige Eenhed af uendelig
Reflexion og barnlig Umiddelbarhed, den vidunderlige
% ---- printed p.26 (PDF 39) ----
Gave, at kjende det Oversandselige paa det Sandselige, og sætte det
Naturlige sammen med det Overnaturlige. Troen forbinder Extremer af
Helligt og Syndigt, af guddommelig Almagt og menneskelig Afmagt: ved
Troen, og kun ved Troen, sættes Mennesket i et paa een Gang endeligt
og dog uendeligt Forhold til Gud som Aand til Aand.

Vi lære det af den hellige Historie. Ved umiddelbar Aabenbaring
faaer Abraham Befaling til at offre sin Søn Isaak. Kunde en saadan
Befaling nu virkelig opfattes med samme umiddelbare, sandselige
Vished, som man kan opfatte en menneskelig Befaling, saa var den
hele Fortælling om Abrahams Lydighed uden religiøs Betydning. Hvad
Abraham gjør, kunde og vilde enhver Anden i hans Sted have gjort.
Thi hvem er det vel, der befaler? Jehova selv, den almægtige Skaber!
Her kan da, under Forudsætning af sandselig Vished, ligesaa lidt
være Tanke om Mulighed for at unddrage sig Befalingen som om Fare
ved at opfylde den. Hvad vovede vel den Mand, der vidste med sig
selv, at den talende Gud var hans Ven, naar han, vel at mærke,
vidste det med sandselig ubedragelig Vished. Den Almægtige kunde jo
% sic: "døer" (for "døe") is the printed reading — verified at 400 dpi, p.26.
maaskee blot ville lade Isaak døer for at vække ham op igjen, eller,
hvis det ikke behagede ham at opvække den Døde, give sin Udvalgte
tusinde Sønner i Erstatning for den Ene. Nei, den religiøse
Betydning ligger netop deri, at Abraham ifølge Sagens Natur umulig
kunde have nogen sandselig afgjørende Vished, og just af den Grund
maatte have Forvisningen i Troen. Han troede at høre Jehovas Stemme;
han troede, at Befalingen, skjøndt den, utroligt nok, løb paa Isaaks
Opoffrelse, udgik fra Gud, og ikke, hvad det naturlige Hjerte dog
maatte finde langt troligere, fra mørke Griller eller fra en ond
Aand. --- Vende vi os til Evangeliet, da behøvede de Samtidige jo
blot at bruge de naturlige Sandser for at see og høre
% ---- printed p.27 (PDF 40) ----
Mennesket Jesus; men til i Menneskens Søn at opdage \emph{Christus,
den levende Guds Søn}, var den naturlige Seen og Høren ikke
tilstrækkelig: \emph{salig er du Simon, Jonas Søn, thi Kjød og Blod
haver ikke aabenbaret dig det, men min Fader, som er i Himmelen}
(Mtth.\ 16).

Saa lidt som en ydre Aabenbaring kan faae religiøs Betydning uden
Tro, ligesaa lidt kan Troen faae religiøs Betydning uden en
tilsvarende indre Aabenbaring. Mangler den indre Aabenbaring, ville
Beretninger om det Overnaturlige enten afvises med Vantro eller paa
overtroisk Maade antages ifølge en blot historisk Tro. Thi al
ligefrem historisk Tro paa den hellige Historie er Overtro. Den
hellige Historie er nemlig deri forskjellig fra anden Historie, at
den vil troes, ikke med, men imod al Sandsynlighed. Paa Grund af
Usandsynligheden er det i Aabenbaringen Historiske, hvad Antagelsen
angaaer, ikke ligefrem understøttende, men tvertimod frastødende.
Det Høieste, der ved historisk Vidnesbyrd kan opnaaes, er kun
Sandsynlighed for, at Vidnerne selv have troet paa det
Usandsynlige. Det Usandsynlige, som ellers er et Kjendetegn paa det
Uhistoriske, er i den hellige Historie netop Mærket paa det
Religiøse, forsaavidt det er Underets Mærke.

De evangeliske Beretninger om Christi Opstandelse gjøre ingenlunde
Opstandelsen selv sandsynlig, saa lidt som de give nogen fuldkommen
Forvisning om dens Virkelighed. Dette er ingen Indvending imod
Beretningernes Troværdighed; tvertimod, de kunne ikke være
troværdigere end de virkelig ere, men Troværdigheden boer i det
Frastødende. Uden de historiske Vidnesbyrd vilde den religiøse Tro
mangle de overnaturlige, ved bestemte Kjendsgjerninger bestemte
Forudsætninger, og dermed sit egentlige Grundlag; men derfor er det
ligefuldt sandt, at det ikke er de udvortes i sig usandsynlige Be\-
% ---- printed p.28 (PDF 41) ----
givenheder, men det i Begivenhederne Indvortes, det Aabenbarede, som
paany aabenbarer sig for Troens Blik, ophæver Forbigangenheden og
gjør det Guddommelige nærværende. Den ydre Aabenbaring er paa
Betingelse af Tro udtrykkelig sat i hellig Historie, for at det
Historiske igjen paa Betingelse af Tro kan omsættes i indre
Aabenbaring.

\runhead{b) Aabenbaring og Fornuft.}

Modsætningen mellem Aabenbaring og Fornuft betegnes kort og skarpt
som en Modsætning mellem ubetinget Frihed og absolut Nødvendighed,
mellem Sandhed i Tro og Sandhed i Viden. At ville tilhylle det
Stridige og Stridbare i denne Modsætning er forgjæves. Man har
beraabt sig paa, at Sandheden maa være een og udeelt, i alle
Henseender overeensstemmende med sig selv; men saasnart man skal til
at angive, hvor det høieste Sandhedskriterium er, i Aabenbaringen
eller i Fornuften, bryder Striden ud.

Naar Spørgsmaalet om Aabenbaringens Sandhed gjøres afhængigt af
Fornuftens Dom, da kommer Videnskaben frem med Indsigelser, den ene
mere høirøstet end den anden. Aabenbaring i Religionens Forstand er
et utænkeligt Begreb, lige selvmodsigende, enten vi see paa
Menneskets eller paa Verdens eller paa Guds Væsen.

„Tage vi“, siger en dybsindig Tænker, „blot Hensyn til den
menneskelige Natur i og for sig, saa overbevise vi os om, at
Aabenbaringen forudsætter en Receptivitet i Menneskets Aand, der
ganske strider mod Aandens Væsen. Thi da Aandens Væsen bestaaer i
Activitet, saa kan den kun være i Liden, forsaavidt den tillige er i
Virken. Men ved en Aabenbaring, en umiddelbar Indvirkning af det
høieste Væsen paa den menneskelige Aand, bliver der Intet tilovers
for den sidste uden en absolut Liden; thi det høieste Væsen er
absolut activt,
% ---- printed p.29 (PDF 42) ----
men det nødvendige Correlat til den absolute Activitet er den
absolute Passivitet.“ (Schelling). Betragtet fra denne Side er en
Aabenbaring altsaa umulig. Ligesaa umulig viser Aabenbaringen sig,
naar vi betragte Verden, den naturlige Tingenes Orden.
Aabenbaringen er et Under, men Underets Begreb ophæver
Naturbegrebet. Her behøves ingen smaalig Kritik til at udfinde
Forskjellen mellem rimelige og urimelige, nødvendige og overflødige
Undergjerninger; Sagen er, at et eneste Under tilintetgjør hele
Naturens Realitet. Thi hvad er Naturen? En fornuftnødvendig
Idee-Aabenbarelse; hvis ikke, da et sandseligt phantastisk
Skyggevæsen. Have Tingene en naturlig Grund og nødvendig
Sammenhæng, saa er Miraklet umuligt; er Miraklet muligt, saa er
Naturgrunden huul, Sammenhængen løs, og Alfornuften en taaget
Indbildning. Et Mirakel, som f.\ Ex.\ det, at en Stav forvandles til
en Slange og Slangen igjen til Stav, er nok til at vise, hvorledes
en naturløs, ubetinget Villie kan lege med alle Betingelser,
behandle de endelige Aarsager som intetsigende Anledningsaarsager,
ophæve Loven og berøve Naturen dens Realitet.

Det samme Urimelige viser sig med Hensyn paa Guds eget Væsen.
Religionen lærer vel, at Gud er uforanderlig; men antage vi dens
Aabenbaringshistorier, da maa Gud i sig selv være meget
foranderlig. En Sammenligning mellem Nutid og Fortid viser noksom,
at der i Tidernes Løb maa være foregaaet en mærkværdig Forandring
med hans Forhold til Verden og Menneskene. Den samme Gud, der fordum
talede og talede næsten uophørlig og gjorde saa mange Mirakler, at
det Gamle og Nye Testamente neppe kunne rumme dem, see, han er i den
senere Tid bleven ganske taus, og gjør slet ingen Mirakler. Fra
Aabenbaringens Standpunkt er dette uforklarligt. Dersom det ikke er
Øiemedet med Miraklerne, at de skulle styrke Troen og bekæmpe
Vantroen,
% ---- printed p.30 (PDF 43) ----
hvad er da Øiemedet? Men dersom en Bestyrkelse i Troen er Øiemedet:
hvo forklarer os da, hvi Miraklerne udeblive? Gaves der nogensinde
en Tid, da Menneskene ret kunde trænge til at høre et virkeligt Guds
Ord, et almægtigt Ord af Guds egen Mund og see en virkelig Guds
Gjerning, da maatte det vist være denne nærværende Tid; thi ligesom
Jordbunden efter langvarig Tørke og Hede trænger til husvalende og
befrugtende Regn, saaledes maae de fromme Gemytter i Tvivlens og
Culturens Solbrand omsider trænge til husvalende og befrugtende
Mirakler. Skal Troen paa en overnaturlig Aabenbaring holdes i Live,
maa den uforanderlige Gud fra Tid til Tid i mærkelig Grad forandre
sig.

Men er en Aabenbaring nu ogsaa nødvendig? Religionen vil begrunde
Nødvendigheden paa Menneskeslægtens almindelige Fordærvelse ved
Syndefaldet; men her, om ellers nogensteds, viser sig et
paafaldende Misforhold mellem Middel og Øiemed. Fordærvelsen, hedder
det, er almindelig, Trangen til den foregivne Hjælp følgelig ogsaa
almindelig; og dog, uagtet Alle ere ved Faldet blevne lige skyldige,
lige hjælpeløse, lige trængende, har Hjælpen dog langtfra udstrakt
sig til Alle. Det var oprindelig kun Jøderne, hvilke Gud kom til
Hjælp ved særegne Aabenbaringer; Hedningene, det store Fleertal,
maatte hjælpe sig selv, saa godt de kunde, ved Fornuftens Brug. Men
at hjælpe Nogle og forbigaae Andre, hvad er det Andet end
guddommelig Partiskhed, grundløs Vilkaarlighed?

At dette i det Hele er „en fornuftig Tankegang“, vil vist Ingen
negte: Fornuften kan som Fornuft, naar den skal gjøre sit eget
Erkjendelsesprincip gjældende, ikke dømme anderledes. Men naar
Fornuften ikke kan dømme anderledes, hvad hjælper det saa at ville
lægge Dølgsmaal paa, at Aabenbaringen er og maa være Fornuften en
Forargelse.
% ---- printed p.31 (PDF 44) ----
Gjøres omvendt Sagen afhængig af Aabenbaringens Kjendelse, da vil
Dommen ligefrem gaae Fornuften imod, og det vil let kunne synes, som
om dens Ret blev krænket. Fornuften mener at kunne forklare, begribe
og gjennemskue Aanden; det er denne Mening, Religionen tilbageviser
som en Indbildning. Thi Aabenbaringen alene forklarer Aanden og
oplader dens Dybder: Aabenbaringen er Aandens Sol. Ligesom det
sandselige Øie blendes, naar det stirrer ind i Solen, saaledes
blendes Fornuften, naar den vil stirre paa Aabenbaringen. Fornuften
veed ikke, hvad Aand er, og kan ikke vide det; thi den veed ikke,
hvad en hellig, almægtig Villie er, og kan ikke vide det. At
Fornuften i sit System af almeengyldige Begreber kan optage ogsaa
Villiesbegrebet, skal ikke negtes; den kan ved Forgudelse af den til
Ideenødvendigheden bundne Almeenvillie endogsaa mene, at den i
Begrebet virkelig anerkjender Guds Villie. Men den Anerkjendelse,
Fornuften yder Gudsvillien ved at berøve den Selvbestemmelsens
Ubetingethed og ophøie den til Rang med „virkelige Begreber“, har
ikke Mere at betyde end den Anerkjendelse, en hedensk Keiser fordum
viste Christendommen ved at optage Christus blandt sine Huusguder.
Det nytter ikke at ville stille den efter Nødvendighedsgrunde
forskende Fornuft tilfreds ved Forsikkringer om, at Guds Villie er
Eet med Guds Væsen, og Guds Væsen og Villie Eet med den evige
Sandhed: hver Gang Villien vil udtrykke sig i \emph{dette} eller
\emph{hiint} Bud, i \emph{denne} eller \emph{hiin} Gjerning, altsaa
som Villie, blive Nødvendighedsgrundene borte; Betingelserne
forsvinde i det Ubetingede; Ubetingethedens Braad saarer Fornuften,
og den saarede Fornuft beklager sig over Vilkaarligheden. Men naar
Fornuften ikke kan begribe den almægtige Villies ubetingede Frihed,
altsaa ikke begribe Aabenbaringens høieste Forudsætning: hvorledes
skulde den da
% ---- printed p.32 (PDF 45) ----
være istand til at magte Religionen, trænge ind i Gudsforholdet og
begribe Troen.

Kun paa een Betingelse kan Guds ubetingede Villie erkjendes; men
denne Betingelse, skjøndt i Grunden een, er dobbeltsidig: fra den
guddommelige Side Aabenbaring, fra den menneskelige Side Tro. Det i
Aabenbaringen afgjørende Første er ikke, hvad Gud i dette eller
hiint Anliggende vil Mennesket, men at Gud udtrykkelig vil Mennesket
Noget, har Villie til Mennesket, altsaa ubetinget: at Gud har
Villie. Aabenbaringen er den guddommelige Frihedsact, den ideale
Handling, hvori Aanden oprindelig forkynder sig som Aand. Troen er,
i Lighed med Aabenbaringen, ikke en uvilkaarlig Rørelse eller uklar
Følelse, men en indre Handling, en første Frihedsact, ved hvilken
Aandsvillien bliver til i Mennesket, og ved hvilken Mennesket
væsentlig bliver Aand. Det i Troen afgjørende Første er ikke dette:
at Mennesket angaaende Dette eller Hiint umiddelbart forvisser sig
om, hvad Gud vil, men dette: at Mennesket skal ville Guds Villie,
ubetinget altsaa: at Mennesket skal \emph{ville}, skal erkjende, at
Gudsforholdet er et Villiesforhold, det Forhold, hvori der alene er
Aand og Frihed og evigt Liv.

I dette Forhold kan Fornuften umulig trænge ind; det seer ud, som
havde det skilt sig fra alle Væsensgrunde: just fordi det optager og
indeslutter alle objective Grunde i Subjectiviteten selv, i
Frihedsgrunden, seer det i Fornuftens Øine ud, som var det
grundløst. Eller hvorpaa skulde vel Troen begrundes? Paa
Aabenbaringen, altsaa paa et Under. Og hvorpaa skulde saa
Aabenbaringens Sandhed begrundes? Paa Troen, altsaa paa et Under,
der selv er begrundet paa et Under. Som ubetinget Villiesforhold er
Gudsforholdet, naar det gribes i sin inderste Eiendommelighed, et
antirationelt Forhold; ved at ville rationalisere det i sig selv
Antirationelle, forraader Fornuften kun sin egen Blindhed. I Aandens
% ---- printed p.33 (PDF 46) ----
% n.b. the lower half of this page LOOKS italic in a downsampled two-up render
% (gutter curvature on a left-hand page); at 420 dpi it is plain roman. No italics.
høieste Anliggender, for alle absolute Villiesafgjørelser er
Fornuften blind, blind for Faldet og blind for Opreisningen; blind
for Faldet, thi Faldets Hemmelighed ligger i Villien; blind for
Opreisningen, thi Opreisningens Hemmelighed ligger i Villien.
Fornuften vil udgrunde Viisdommen, og kan dog med al sin Viisdom
ikke erkjende Gud i hans Viisdom; thi alle Viisdommens Hemmeligheder
skjule sig i Villien, Viisdomsideen, Guds Villies Idee, aabenbares
kun for Villien. Saa dømmer Religionen „den blinde Fornuft“. Det vil
fra dette Synspunkt kunne forstaaes, at det ikke er Fornuften, men
dens Overgreb, der er Aabenbaringstroen en Forargelse.

De mange vidtløftige og tilsyneladende endeløse Stridigheder mellem
Aabenbaring og Fornuft kunne kun da bringes til en sidste
Afgjørelse, naar de føres tilbage til det første Udgangspunkt,
Modsætningen nemlig mellem Frihedens og Nødvendighedens Principer,
mellem Villiesgrunde og Væsensgrunde. Enhver afgjørende Villiesgrund
trækker altid det Almene ind i Enkeltheden, hæver Nødvendigheden ved
det frie Valg og gjør Væsensgrunden utilstrækkelig; i den
tilstrækkelige Væsensgrund er det Enkelte derimod opløst i det
Almene, Nødvendigheden sat med de givne Betingelser, og
Villiesgrunden ukjendelig. Begge Dele kan Fornuften indsee; men
derfor kan den ogsaa indsee, at den i Villien ophævede Grund, der
giver Handlingen som Handling Præg af Vilkaarlighed, umulig kan være
Fornuftgrund. I denne Sammenhæng vil det tillige kunne indsees, hvor
ørkesløs og tom den Tænkning er, som bilder sig ind at kunne hæve
Skilsmissen imellem de ueensartede Principer ved den Forsikkring, at
Frihed og Nødvendighed absolut maa være Eet i det absolute Væsen.
Thi saasnart der spørges, hvad det da egentlig er, som gjør det
absolute Væsen absolut, saa
% ---- printed p.34 (PDF 47) ----
sætter, som allerede viist, Aabenbaringen igjen Friheden, Fornuften
Nødvendigheden, og Striden begynder forfra.

Her er ingen anden Udvei end at lade de to absolut ueensartede
Principer afgive Midtpunkter for to absolut ueensartede
Begrundelser, to absolut ueensartede Erkjendelseskredse. Er det nu
ligesaa umuligt at aflede Nødvendigheden af Friheden som Friheden af
Nødvendigheden, og tillige umuligt at negte de modsatte Principers
oprindelige Grundeenhed: saa maa Grundeenheden nødvendigviis
opfattes som et for den menneskelige Viden uigjennemtrængeligt
Mysterium.

Dette Mysterium er intet Paradox. Thi Paradoxet viser sig kun, naar
den vidende Bevidsthed eensidig tager Parti for Fornuften, og vil,
ved at anstrenge Forstanden til det Yderste, af Fornuftprincipet
begribe Aabenbaringen. Ved at anerkjende Mysteriet anerkjende vi
begge Erkjendelseskredsenes indre Uendelighed. Alle af
Nødvendighedsprincipet udviklede Erkjendelser have i deres egen
Sphære fuld Gyldighed, og Fornuften forsaavidt en ubestridelig Ret
til at undersøge Alt, prøve Alt og udvikle sin Videns Indhold i det
Uendelige. Ligeledes have alle af Frihedsprincipet udviklede
Erkjendelser fuld Gyldighed i deres Sphære, og Aabenbaringstroen fra
sin Side en ubestridelig Ret til at lade sin Tænkning prøve Alt og
dømme om Alt; Fornuften kan i denne Henseende ikke sætte den
conseqvente Udvikling af Troens Indhold nogensomhelst Skranke.

Men, spørger man, naar Aabenbaring og Fornuft nu komme i Strid om
det samme, det selvsamme Virkelige, de samme, de selvsamme
Kjendsgjerninger: hvad saa? Ja, saa følger heraf, at det samme
Virkelige maa vise sig forskjelligt og bedømmes forskjelligt og
tillægges forskjellig Betydning, eftersom det sees enten i
Aabenbaringens eller i Fornuftens Lys. At Aabenbaring og Fornuft
modsige hinanden, er da saa langt fra at være et Beviis
% ---- printed p.35 (PDF 48) ----
enten imod Aabenbaringens religiøse eller imod Fornuftens
videnskabelige Sandhedserkjendelse, at det meget mere er en
nødvendig Følge af Principernes absolute Ueensartethed.

% sic: the run-head prints "Aabenbaringen" without the genitive -s; the forward
% reference on p.24 reads "Aabenbaringens Tilegnelse i Troen". Verified at 420 dpi.
\runhead{c) Aabenbaringen Tilegnelse i Troen.}

Hvor ganske nøiagtig Aabenbaring og Tro svare til hinanden, sees
udtrykkelig deraf, at Begreberne selv i deres Gjensidighed defineres
ved hinanden. Hvad er vel det Aabenbarede? Just det, som alene kan
opfattes af Troen, erkjendes, bevares og forklares i Troen, det,
hvis Sandhed alene kan svare til Troens Sandhed. Og Troen? Ja,
ligesom Synet særlig bestemmes ved det Synlige, Hørelsen ved det
Hørlige, Tanken ved det Tænkelige, saaledes maa Troen bestemmes ved,
hvad den naturlige Sands anseer for det Utrolige, omendskjøndt det
dog ifølge sin Natur netop er det i Sandhed Trolige, det Eneste, som
i Ordets egentligste og strengeste Forstand lader sig troe, nemlig
det af Gud selv Aabenbarede. Ved denne Grundbestemmelse er ethvert
Forsøg paa at indblande Fornuften i Religionens Anliggender een Gang
for alle afskaaret. Men herved er tillige stillet
Religionsphilosophien den Fordring, at den skal vise, hvorledes
Forholdet imellem Aabenbaring og Tro lader sig udvikle af
Gudsforholdet selv, og derved aabne os et Indblik i Troens
Psychologie.

Betragte vi de to Hovedsider i Gudsforholdet, den guddommelige og
den menneskelige, hver for sig, da tyder den første særlig hen paa
det Udvortes, den sidste paa det Indvortes i Aabenbaringsunderet;
Modsætningen imellem det Udvortes og det Indvortes er her saa
fremtrædende, at den begrunder en Adskillelse mellem det aabenbare
og det skjulte Under.

Gudsforholdet sættes ved en guddommelig Handling; det Evige tager
sin Begyndelse i Tiden, det Over\-
% ---- printed p.36 (PDF 49) ----
historiske bliver historisk. Dette er Fornuften en forargelig
Modsigelse, og Modsigelsen skærpes til det Yderste, naar Blikket
fæster sig paa en Kjendsgjerning som den, at Gud til en given Tid
har aabenbaret sig i et enkelt Menneske, et Menneske i Kjød og Blod;
dette Menneske, der er født, har lidt og er død, skal virkelig være
Gud! det seer jo ud, som var Modsigelsen selv bleven incarneret. Er
det Evige paa denne Maade blevet historisk, saa er det, siger den af
Fornuften ophidsede Forstand, blevet til Noget, som udtrykkelig
strider imod dets Natur, og maa altsaa være blevet det „i Kraft af
det Absurde“. --- Men det i sig selv Absurde, er \textit{eo ipso}
det i sig selv Umulige, og det i sig selv Umulige kan ikke troes.
Her maa da vise sig en Forskjel imellem det i Fornuftens Tanker
Umulige og det i sig selv Umulige. Skal den Eviges Selvaabenbarelse
i Tiden virkelig høre til det i sig selv Umulige, da maae saadanne
religiøse Udtryk som: guddommelig Gjerning, guddommelig Handling,
guddommelig Villie nødvendigviis henregnes til det i sig selv
Umulige. Eet af To: enten er Gud Villie, og da kan han ogsaa
incarnere sig; eller Gud kan ikke incarnere sig, og da er han heller
ikke Villie. Underet som Phænomen er et Udvortes, men Phænomenets
Aarsag, Villien selv, er et Indvortes. Forargelsen paa det Udvortes,
paa det aabenbare Under, er altsaa en Forargelse paa det indvortes,
det skjulte Under, en Forargelse paa den absolut religiøse
Forudsætning, at den almægtige Villie er Gud, at Guds Villie er fri,
at Gud kan handle.

I det ved Aabenbaring og Tro bestemte Gudsforhold er den frie,
almægtige Handling, at Gud træder Mennesket imøde, ubetinget det
Første. I Synet, i Tegn og Tale skeer Meddelelsen umiddelbart. Dette
Syn, dette Tegn, dette Ord henvendes ikke til Mennesker i
Almindelighed, men til dette bestemte Menneske, angaaer ikke
menneskelige Anliggender i Almindelighed, men dette
% ---- printed p.37 (PDF 50) ----
eller disse bestemte Anliggender: Aabenbaringen kommer fra
Uendeligheden som et Lyn, der slaaer ned i det Endelige. Synet,
Tegnet, Ordet, dette Overnaturlige skeer ved umiddelbart
Gjennembrud af det Naturlige; Underet er aabenbart, og Eftertrykket
falder paa det Udvortes. Den Tro, der umiddelbart tilegner sig en
saadan udvortes Aabenbaring, maa ifølge den rationelle Psychologie
conseqvent opfattes som Overtro. Skulde Muligheden af en udvortes
Aabenbaring regnes med til den fornuftige Tingenes Orden, da maatte,
som ovenfor anført, slige Aabenbaringer jo ligesaa vel kunne skee
til en Tid som til en anden, ligesaa vel blive det ene Menneske til
Deel som det andet; men Enhver, der i Livets virkelige Anliggender
beraaber sig paa Syner og Aabenbaringer, bliver anseet, om ikke
ligefrem for en forsætlig Bedrager, saa dog for en overtroisk
Sværmer.

Vende vi os nu til det Næste, nemlig Gudsforholdets menneskelige
Side, faaer Sagen et modsat Udseende. Det Udvortes forsvinder med
dets overordentlige Indvirkninger, Gudsvillien forvandler sig til en
ved det Godes Idee bestemmelig og dog overnaturlig Almeenvillie; Gud
drager sig tilbage; Alt foregaaer i den Troendes Indre: det
aabenbare Under giver Plads for det skjulte. Fra den rationelle
Psychologies Standpunkt er Dommen over en saadan Sjæletilstand
betegnende udtalt i Strausses Kritik af Schleiermachers Tro paa den
historiske Christus. „Jeg føler som Christen“, siger Schleiermacher,
„Foreningen af min lavere og min høiere Selvbevidsthed lettet, eller
min Fromhed befordret. Denne Befordring kan ikke udgaae fra mig
selv, da der tvertimod fra mig kun udgaaer en Standsning af det
religiøse Liv“. Strauss spørger: „Hvorfra veed Du det? Du nemlig,
det moderne-christelige Jeg, thi det gammeldagstroende vidste
rigtignok sammestedsfra som det vidste, at det af sig selv ikke
vidste noget Sandt, at det af sig selv ikke kunde virke
% ---- printed p.38 (PDF 51) ----
noget Godt: begge Dele vidste det af sin Grundforudsætning, Aandens
Fremmedhed for sig og Udtræden fra sig. Men hvorfra faaer det
moderne Jeg, der er kommet til Bevidsthed om sig og er blevet sig
mægtigt, Midler til at foretage en saadan Adskillelse? Om noget Ydre
og Tilfældigt, som en Efterretning, en historisk Notits, kan jeg vel
vide, at jeg ikke har den fra mig selv, men har erfaret den af
Andre; men her handles der om to modsatte Bevægelser i det
menneskelige Gemyts Inderste, af hvilke kun den ene skal have sin
Oprindelse der, den anden derimod være ledet derind udenfra“. Det
religiøse Sindelag bliver altsaa kjendeligt paa det
„Rettroenhedens barnagtige Regnestykke, ved hvilket Mennesket først
\textit{per subtractionem} skiller sig af med alt Godt, for derpaa
\textit{per additionem} at lade det tilflyde sig som Naadegave“. Det
skjulte Under er den uklare Bevidstheds egen skjulte Modsigelse: en
saadan Reflexion maa conseqvent opløse Troen og ende i Vantro.

Ved de Extremer, som i Conseqvens af Gudsforholdets Natur maae danne
sig mellem det Ydre og det Indre, bestemmes Troen. Det Ydre,
Gudsindvirkningen i dens umiddelbare Særskilthed, er nødvendig
Betingelse for at vække det Indre; men det er som Ydre kun
Betingelse, kun Middel; det Indre og dets Tilegnelse er Maalet.
Synet er et Ydre, men Almagtens Væsen et Indre; Ordet, der høres, er
et Ydre, men Tanken i Ordet, den guddommelige Tanke, et Indre;
Befalingen et Ydre, men Villien, der befaler, et Indre. Troen er nu
deri forskjellig fra Overtroen, at den igjennem det Ydre griber og
tilegner sig det Indre; og deri forskjellig fra Viden, at det Indre,
den griber, er Overnaturligt, altsaa et Indre, der ikke ligefrem
eller i Kraft af fornuftnødvendig Reflexion bestemmes ved sit Ydre:
skulde Troen indlade sig paa objective Væsensreflexioner, maatte den
ende i Vantro.
% ---- printed p.39 (PDF 52) ----
Idet Overtro og Vantro begge falde udenfor Troen, fremkalde de
hinanden gjensidig. Overtroen vil have sandselig Vished om det
Overnaturlige, et Tegn for hvert Tilfælde, et Varsel for hver
Begivenhed, en Guds Befaling for hver Handling; den hænger sig ved
det Udvortes. Forstand og Fornuft blive da let enige om, at saadanne
Forestillinger om det Overnaturlige ere unaturlige og vanære det
Guddommelige. Saa forkastes Tegn og Varsler, og de sædelige
Handlinger henføres under sædelige Love. Med de sædelige Love følger
Ideenødvendigheden; men Ideenødvendigheden er som en Sky, der
skjuler Guds Aasyn: „Gud har ingen Villie, der er ingen Gud“, det er
Vantroens Forklaring, og denne Forklaring er tom og trøstesløs. Det
tomme, gudforladte Sind griber da atter efter overnaturlige
Forudsætninger, vil atter have sandselig Vished: Vantro føder igjen
Overtro, og saaledes i det Uendelige. Af Troen som Tro kan Overtroen
derimod aldrig fødes; thi Troen omsætter det Udvortes i sit
tilsvarende Indvortes og lever af Aanden; heller ikke kan Vantroen
nogensinde faae Magt med Troen; thi Troen har Inderlighed, og
Inderligheden hviler paa overnaturlige Kjendsgjerninger, paa frie,
guddommelige Forudsætninger, der modstaae Nødvendigheden.

Er nu det Indre sat som det oprindelig Første, men dog saaledes, at
det tillige er betinget af et ophævet Ydre, saa maa der jo i et
saadant Indre kunne paavises et Sandhedskriterium, et ubedrageligt
Mærke paa Sandt og Falskt; dette Kriterium maa være bestemt ved det
til Aabenbaringsprincipet svarende Forhold mellem Indre og Ydre.
Hvori skal Sandhedskriteriet da søges? Ikke i Følelsen; thi naar
Følelse strider mod Følelse, alt som Stemningerne vexle, saa behøves
her just et Kriterium for Følelsen%
% This footnote begins on printed p.39 and runs over onto p.40 (PDF 52--53).
\footnote{At Religionens Væsen ikke udelukkende kan sættes i
Følelsen, selv om vi tage Følelsen i substantiel Almindelighed som
absolut Afhængighedsfølelse, har Martensen tydelig viist. „Naar
Schleiermacher har betegnet den religiøse Følelse som en Følelse af
absolut Afhængighed, da har han derved ligesom Mystikerne betegnet
Fromheden som en gudlidende Tilstand, en Tilstand, hvori Mennesket i
sit Gemyts Inderste føler sig berørt af den Magt, i hvilken vi leve,
røres og ere, en hellig Pathos, hvori Mennesket føler sig som
Guddommens Kar og Bolig. Men som denne Beskrivelse minder om
Mystiken, saaledes er den selv mystisk; thi det staaer i ubestemt
Dæmring, \emph{hvilken} denne absolute Magt er, af hvilken jeg føler
mig afhængig, om det er det upersonlige Absolute, om det er Fatum,
eller om det er den ethiske, den hellige, den gode Magt. Kun i
sidste Tilfælde kan den gudlidende Tilstand, kan den absolute
Afhængighedsfølelse tillige være en befriende og opløftende
Følelse.“ Den christelige Dogmatik. Indledning S.\ 10.\par
Hvor vanskeligt det falder Schleiermacher at gjøre Adskillelsen
mellem Guds og Verdens Væsen ret kjendelig, fremlyser da ogsaa af,
hvad han („\textit{Der christliche Glaube}“. Berlin 1835. 1.\ Bd.\
S.\ 169) siger om „den deelte og den udeelte Eenhed“; thi vel er
Verdenseenheden deelt, men, seet fra Synspunktet af det immanent
Absolute, er den dog ogsaa udeelt, og Umuligheden af en absolut
Afhængighedsfølelse i Forhold til Verden saaledes ingenlunde
beviist.}.
Endnu mindre i Phantasien, da Phanta\-
% ---- printed p.40 (PDF 53) ----
% (the body text of p.40 is only four lines; the rest of the page carries the
%  continuation of the footnote begun on p.39, transcribed in full above)
sien ligesaa let kan bevæges af usande som af sande Indskydelser. At
søge Kriteriet i den almene Tænkning hjælper heller ikke, eftersom
Tankens Bevægelser kunne være ligesaa forskjellige som de Principer,
hvoraf den bevæges. Kriteriet maa, ifølge Forholdets Natur,
oprindelig søges i Villien.

Som almindelig Villiesact kan Troen kun bestemmes fra Grunden af ved
Villien selv, nemlig ved Grundvillien. I Grundvillien bliver det
Inderste i Mennesket, det menneskelige Selv, umiddelbart berørt,
beaandet og bevæget af det Inderste i Gud, af det guddommelige Selv,
af Gudsvillien.

Men naar Adskillelsen af Sandt og Falskt gjøres afhængig af Villien,
maa den da ikke blive vilkaarlig? Ingenlunde! Den Villie, hvorom her
tales, er ikke den egoistiske, mellem endelige driftagtige
Tilskyndelser væl\-
% ---- printed p.41 (PDF 54) ----
gende Formalvillie, men Aandsvillien selv, den alle Sindets
Kræfter, alle Følelser, Forestillinger og Tanker gjennemtrængende,
det hele sjælelige Liv samlende Frihedsvillie, den Villie, i Kraft
af hvilken det indvortes Menneske ikkun kan og maa ville sin egen
Frelse. Vækkelsen af denne Aandsvillie er udelukkende betinget af
Gudsvilliens Indvirkning; den beroer altsaa paa et skjult Under;
heraf sees dens Forhold til historiske Kjendsgjerninger af aabenbare
Undere; den er et Indre, der i Troens Forstand reflecterer sit
Ydre. Aandvilliens Eenhed med Følelse, Phantasie og Tanke giver i
Ordets allerdybeste Betydning Selvforstaaelse, en Selvforstaaelse,
der alene beroer paa inderlig Forstaaelse med Gud. Til videnskabelig
Prøvelse af Sandt og Falskt opstiller Fornuftphilosophien
\emph{Begrebsmodsigelsens Grundsætning}; til personlig Prøvelse af
Sandt og Falskt --- til Selvprøvelse for Gud --- opstiller
Religionsphilosophien \emph{Selvmodsigelsens
Grundsætning}\footnote{Ved Begrebsmodsigelser tænkes her ikke blot
paa Modsigelser mellem rene objective Vidensbegreber saasom: en rund
Fiirkant, et udenfor Rummet faldende Legeme o.\ s.\ v., men tillige,
og det fornemmelig, paa saadanne Modsigelser, som fremkomme ved en
Sammenblanding af Religion og Videnskab, idet man (jvnfr.\
Forelæsningerne: Om Hindringer og Betingelser for det aandelige Liv
i Nutiden) af religiøse Forudsætninger uddrager videnskabelige, og
af videnskabelige Forudsætninger religiøse Conseqvenser. Istedenfor
ægte Vidensbegreber og rene Troesbegreber faaer man da en Række
Bastardbegreber, der naturligviis alle ere fulde af Modsigelser. At
en videnskabelig Modsigelse iøvrigt, naar man tager det formelt,
ogsaa kan opfattes som en Selvmodsigelse, og omvendt enhver
Selvmodsigelse som Begrebsmodsigelse, forstaaer sig af sig selv, og
bliver kun udtrykkelig anført for at spare „Ordkritiken“ den
Uleilighed at gjøre formelle Opdagelser, hvor Opdagelser ere
overflødige.}.

I Selvforstaaelse for Gud er Frelse, i Selvmodsigelse for Gud er
Fortabelse; det er Aandslivets yderste Poler; om disse to Poler maae
i Troens Inderlighed alle Følelser,
% ---- printed p.42 (PDF 55) ----
al Frygt, alt Haab, alle Hjertets Tanker bevæge sig. Det
Eiendommelige ved Gudsforholdet viser sig nu deri, at Alt paa een
Gang er saa umiddelbart og dog uendelig reflecteret, saa enfoldigt
og dog, naar Eftertanken ret begynder, uendelig indviklet. Troen er
af barnlig Natur: hvad der er skjult for de Vise og Forstandige, er
aabenbaret for de Umyndige; men Troen er dog først og sidst af
aandelig Natur; dens Væsen er Aand, og „Aanden randsager alle Ting,
ogsaa Guds Dybheder“.

\subhead{Ordet og Aanden.}
\parmark{4}

Men „Guds Dybheder“ i Aanden randsages ikke paa samme Vilkaar som
Ideens Dybheder i Fornuften; thi Fornuftideen er fordybet i
Nødvendighed, men i Gud er Inderligheden fri, uendelig fordybet i
Selvbestemmen. Til et ufrit Indre svarer et ufrit Ydre, og
Yttringen skeer med Nødvendighed; saaledes maa Væsenet nødvendigt
have sit Skin, Realgrunden sin Existens, Sjælen sit Legeme o.\ s.\
v. Derimod kan det frie Indre kun give sig et sandt Udtryk i et frit
Ydre, og Yttringen skeer med Frihed. Det ufrie Indre lader sig under
givne Betingelser tvinge frem til Aabenbarelse i sit Ydre; derfor
kan Naturen, om vi saa maae sige, ingen Hemmeligheder have for den i
sine Forskninger stedse mere indtrængende, altid videre gaaende
Fornuft. Men det frie Indre kan ikke saaledes udforskes. At et
Menneske under visse Omstændigheder lader sig overraske og kommer
imod sin Villie til at forraade sit Indre, viser kun, at dets Frihed
er ufuldkommen; thi Guds Indre, Friheden selv i dens indre
Uendelighed, skal ingen nok saa dybsindig Fornuft bringe til at
forraade sig. Hvad der skal kunne vides om Gud, maa være aabenbaret
af Gud, aabenbaret

\end{document}
